\section{Results}
\label{sec:results}

\subsection{Gas Cost by Layout}

Table~\ref{tab:gas-first-batch} presents the gas measurements for the first batch
submission (cold storage access pattern) across all three layouts.

\begin{table}[htbp]
\centering
\caption{Gas Cost -- First Batch Submission (Cold Storage)}
\label{tab:gas-first-batch}
\begin{tabular}{lrrrcc}
\toprule
\textbf{Layout} & \textbf{Storage} & \textbf{ZK Verify} & \textbf{Total} & \textbf{Bytes} & \textbf{$<$300K} \\
\midrule
A: Minimal   & 80,156  & 205,600 & 285,756 & 32 & Yes \\
B: Rich      & 102,799 & 205,600 & 308,399 & 64 & No \\
C: Events    & 57,887  & 205,600 & 263,487 & 0  & Yes \\
\bottomrule
\end{tabular}
\end{table}

Only Layouts A and C satisfy the 300K gas target on the first batch. Layout B exceeds the
budget by 8,399 gas (2.8\%) due to the additional \texttt{SSTORE} operation for the packed
\texttt{BatchInfo} struct.

Table~\ref{tab:gas-steady-state} presents steady-state gas (10th batch), where enterprise
state slots are warm and only new batch-specific mapping entries incur cold
\texttt{SSTORE} costs.

\begin{table}[htbp]
\centering
\caption{Gas Cost -- Steady State (10th Batch)}
\label{tab:gas-steady-state}
\begin{tabular}{lrrr}
\toprule
\textbf{Layout} & \textbf{Storage+Logic} & \textbf{Total Estimate} & \textbf{$<$300K} \\
\midrule
A: Minimal   & 63,056 & 268,656 & Yes \\
B: Rich      & 88,053 & 293,653 & Yes \\
C: Events    & 40,787 & 246,387 & Yes \\
\bottomrule
\end{tabular}
\end{table}

In steady state, all three layouts satisfy the 300K target. The cold-to-warm transition
reduces Layout A gas by 17,100 (from 80,156 to 63,056), corresponding to the enterprise
state struct transitioning from cold to warm \texttt{SLOAD}/\texttt{SSTORE} access.

\subsection{Gas Decomposition}

Figure~\ref{fig:gas-breakdown} presents the gas decomposition for Layout A (first batch),
revealing the dominant cost component.

\begin{table}[htbp]
\centering
\caption{Gas Breakdown -- Layout A First Batch (285,756 total)}
\label{fig:gas-breakdown}
\begin{tabular}{lrr}
\toprule
\textbf{Component} & \textbf{Gas} & \textbf{\% of Total} \\
\midrule
BN256 pairing (4 pairs)      & 181,000 & 63.3\% \\
BN256 scalar mul (4$\times$)   & 24,000  & 8.4\% \\
BN256 point add (4$\times$)    & 600     & 0.2\% \\
\cmidrule{1-3}
\textit{ZK verification subtotal} & \textit{205,600} & \textit{71.9\%} \\
\midrule
SSTORE batchRoots (0$\to$nz)  & 22,100  & 7.7\% \\
SSTORE currentRoot (nz$\to$nz) & 5,000   & 1.8\% \\
SSTORE batchCount (warm)      & 2,900   & 1.0\% \\
SSTORE totalBatches (warm)    & 2,900   & 1.0\% \\
SLOAD enterprise state (cold) & 4,200   & 1.5\% \\
Event emission                & 2,150   & 0.8\% \\
Calldata + tx base + logic    & 40,906  & 14.3\% \\
\cmidrule{1-3}
\textit{Storage+logic subtotal}   & \textit{80,156}  & \textit{28.1\%} \\
\midrule
\textbf{Total}                & \textbf{285,756} & \textbf{100\%} \\
\bottomrule
\end{tabular}
\end{table}

ZK proof verification accounts for 71.9\% of the total gas. Within the verification cost,
the BN256 pairing check alone (181,000 gas) represents 63.3\% of total gas. The entire
storage and logic budget---including all \texttt{SSTORE} operations, \texttt{SLOAD} reads,
event emission, calldata costs, transaction base cost, and Solidity computation---fits
within the remaining 28.1\%.

\subsection{Delta Analysis}

Table~\ref{tab:delta} isolates the marginal cost of each storage design decision.

\begin{table}[htbp]
\centering
\caption{Marginal Gas Cost of Storage Decisions}
\label{tab:delta}
\begin{tabular}{lrr}
\toprule
\textbf{Comparison} & \textbf{First Batch} & \textbf{Steady State} \\
\midrule
Rich $-$ Minimal (metadata cost)      & +22,643 (+28.2\%) & +24,997 (+39.7\%) \\
Minimal $-$ Events (root history cost) & +22,269           & +22,269 \\
\bottomrule
\end{tabular}
\end{table}

The root history cost delta (+22,269 gas) maps precisely to one \texttt{SSTORE}
zero-to-nonzero operation (22,100 gas base + overhead), confirming that each additional
storage slot per batch costs approximately 22K gas. The metadata cost delta (+22,643 first
batch) includes the \texttt{BatchInfo} struct write plus the additional \texttt{SLOAD} for
the previous batch's cumulative transaction count.

\subsection{Delegated Verification Estimate}

We analytically estimate the cost of delegated verification (cross-contract call to a
separate \texttt{ZKVerifier} contract):

\begin{table}[htbp]
\centering
\caption{Delegated vs. Integrated Verification (Analytical)}
\label{tab:delegated}
\begin{tabular}{lrr}
\toprule
\textbf{Architecture} & \textbf{Est. Gas} & \textbf{$<$300K} \\
\midrule
Integrated (Layout A)     & 285,756  & Yes \\
Delegated (A + ZKVerifier) & $\sim$341,856 & No \\
\midrule
\textit{Delta}            & +56,100  & --- \\
\bottomrule
\end{tabular}
\end{table}

The delegated architecture adds approximately 56K gas: 2,600 gas for the cold
\texttt{CALL} to \texttt{ZKVerifier}, 49,200 gas for three \texttt{SSTORE} operations in
the \texttt{ZKVerifier}'s \texttt{BatchVerification} struct, and 4,300 gas for two
additional events. This overhead exceeds the entire 14,244 gas margin between Layout A
(285,756) and the 300K target.

\subsection{Invariant Verification}

All seven safety invariant tests passed across all three layouts (Table~\ref{tab:invariants}).

\begin{table}[htbp]
\centering
\caption{Invariant Test Results}
\label{tab:invariants}
\begin{tabular}{lccc}
\toprule
\textbf{Invariant} & \textbf{Layout A} & \textbf{Layout B} & \textbf{Layout C} \\
\midrule
ChainContinuity (gap)     & Pass & Pass & Pass \\
NoReversal                & Pass & Pass & Pass \\
EnterpriseIsolation       & Pass & ---  & ---  \\
HistoryQueryability       & Pass & Pass & N/A  \\
EventRecovery             & ---  & ---  & Pass \\
CumulativeTracking        & ---  & Pass & ---  \\
\cmidrule{1-4}
\textit{Total}            & 5/5  & 4/4  & 4/4  \\
\bottomrule
\end{tabular}
\end{table}

The ChainContinuity invariant---$\mathit{prevRoot}_{\text{submitted}} =
\mathit{currentRoot}_{\text{stored}}$---is both necessary and sufficient for gap detection
and reversal prevention. The \texttt{RootChainBroken} custom error reports both the expected
and provided roots, enabling the enterprise node to diagnose and recover from submission
failures.
